\ysrc{15}
，并且是 $X_{\overline{\Q}_\ell}$ 的一个不可约分支。

在上一命题的证明中，我们证明了 $Y_{\overline{\Q}_\ell}$ 含有点 $(F,\sigma)$，其中 $F$ 形如 $F = F_\lambda c' = c'F_\lambda$，这里 $F_\lambda = \lambda(q^{\frac12})$，$\lambda$ 是 $H_{\overline{\Q}_\ell}$ 的某个余特征标（cocharacter），而 $c'$ 的阶有限。把这个点在 $H^0(\overline{\Q}_\ell)$ 下共轭，可以假设 $\lambda$ 通过 $H_{\overline{\Q}_\ell}$ 的一个极大环面分解，而该极大环面来自 $H_{\overline{\Z}_\ell}$ 的一个极大环面。于是 $\lambda$ 本身必须定义在 $\overline{\Z}_\ell$ 上。由于 $q$ 是 $\ell$-进单位，可知由 $F$ 生成的群在 $H(\overline{\Q}_\ell)$ 中关于解析拓扑的闭包 $\langle F\rangle$ 是紧的。另一方面 $\sigma$ 本身是有限阶元素与幂幺元素的乘积 $\sigma = \sigma_{ss}\sigma_u$，因此 $\langle\sigma\rangle$ 紧；最后 $\langle F,\sigma\rangle$ 作为 $\langle F\rangle\ltimes\langle\sigma\rangle$ 的商也紧。特别地，它含于某个 $H(L)$ 中，其中 $L$ 是 $\Q_\ell$ 在 $\overline{\Q}_\ell$ 中的有限子扩张，并且它稳定 $H^0$ 在 $L$ 上之楼的某个面（facet）。也许再把 $L$ 扩大，可以假设它稳定一个超特殊（hyperspecial）顶点，且该顶点在 $H^0(L)$ 下共轭于与整模型 $H_{\overline{\Z}_\ell}$ 相联系的「原点」顶点。但原点顶点的稳定化子是 $H(\OO_L)Z(H)(L)$，其中极大的有界子群为 $H(\OO_L)$。于是在 $H^0(L)$ 下共轭之后，我们得到一点 $(F,\sigma)\in Y(\overline{\Q}_\ell)$ 使 $\langle F,\sigma\rangle\subset H(\overline{\Z}_\ell)$。这正是说 $(F,\sigma)\in Y(\overline{\Z}_\ell)$，如所欲证。\hfill$\square$

至于余下的论断，我们把它留作习题：对诺特整环 $A$，或者 $\ell$ 在 $A$ 中可逆，或者 $A$ 是 $\ell$-进分离的。在我们的情形，$\Z_\ell$-点的存在表明 $\ell$ 在环 $\OO(Y)$ 中不可逆，故后者是 $\ell$-进分离的。由于 $R_{H,1}$ 是约化的，映射
\[
R_{H,1}\longrightarrow \prod_Y \OO(Y)
\]
是单射，故 $R_{H,1}$ 也是 $\ell$-进分离的。
\end{proof}

现在我们要把泛态射 $\varphi_{H,1}:W^0_t\to H(R_{H,1})$ 延拓到非离散化的驯顺 Weil 群 $W_t = W_F/P_F$。

\begin{prop}{3.9}\label{prop:dat-3-9}
固定一个素数 $\ell\neq p$。存在唯一的 $\ell$-进连续延拓 ${}^\ell\varphi_{H,1}:W_t\to H(R_{H,1}\otimes\Z_\ell)$，它是 $\varphi_{H,1}$ 的延拓；并且当 $R$ 取遍 $\Lambda\otimes\Z_\ell$-代数时，它在 $\ell$-进连续态射 $W_t\to H(R)$ 之中是泛的。此外，它对 ind-$\ell$-进拓扑（见 2.2）也是连续且泛的。
\end{prop}

\begin{proof}
令 $\sigma := \varphi_{H,1}(s)\in H(R_{H,1})$。在命题 3.3 的证明过程中，我们证明了存在与 $p$ 互素的数 $M$，使 $\sigma^M$ 逐纤维幂幺，即对 $R_{H,1}$ 上任一代数闭域 $L$，它在 $H(L)$ 中的像幂幺。特别地，对一切 $m\in\N$，$\sigma^M$ 在 $H(R_{H,1}/\ell^mR_{H,1})$ 中的像具有有限的 $\ell$-幂阶，因此 $\varphi_{H,1,\mathrm{mod}\,\ell^m}:W^0_t\to H(R_{H,1}/\ell^mR_{H,1})$ 唯一地延拓到 $W_t$，并且它在 $I_F/P_F = s^{\hat\Z_p}$ 上的限制通过
\[
\hat\Z_p := \widehat{\Z[\tfrac1p]}\longrightarrow \Z/M\Z\times\Z_\ell
\]
分解。取极限，得到态射 ${}^\ell\hat\varphi_{H,1}:W_t\to H((R_{H,1}\otimes\Z_\ell)^\wedge)$，其中 $\wedge$ 表示 $\ell$-进完备化。我们需要说明这个延拓通过 $H(R_{H,1}\otimes\Z_\ell)$ 分解；由于 $R_{H,1}$ 是 $\ell$-进分离的，后者是 $H((R_{H,1}\otimes\Z_\ell)^\wedge)$ 的子群。只需在与 $H$ 在 $\Lambda$ 上的一个忠实表示 $\iota:H\into \mathrm{GL}_N$ 复合之后来证明。但此时，$\iota\circ{}^\ell\hat\varphi_{H,1}$ 的延拓在 $s^z$（$z\in\hat\Z_p$）上由如下公式显式给出
\[
\iota\sigma^z = \iota\sigma^a\,\iota\sigma^{Mz'} = \iota\sigma^a\bigl(1+(\iota\sigma^M-1)\bigr)^{z'} = \iota\sigma^a\sum_{k=0}^{N-1}\binom{z'}{k}_{\!\ell}\,(\iota\sigma^M-1)^k
\]
其中 $z = a+Mz'\in\hat\Z_p$，$a\in\{0,\cdots,M-1\}$，而 $z'_\ell$ 是 $z'$ 在 $\Z_\ell$ 中的像；这里用到 $\iota\sigma^M$ 是幂幺矩阵。注意广义二项式系数 $\binom{z'}{k}_{\ell}$ 落在 $\Z_\ell$ 中，故该公式给出 $\mathrm{GL}_N(R_{H,1}\otimes\Z_\ell)$ 中的一个矩阵，如所欲证。
\end{proof}
\ysrc{16}
${}^\ell\varphi_{H,1}$ 的泛性的验证相当形式化，留给读者。最后，为证明 ${}^\ell\varphi_{H,1}$ 对 ind-$\ell$-进拓扑的连续性，只需说明：对任意向量 $v\in(R_{H,1}\otimes\Z_\ell)^N$，由集合
\[
\{\iota({}^\ell\varphi_{H,1}(s^z))v,\ z\in M\hat\Z_p\}
\]
生成的 $\Z_\ell$-子模是有限生成的。但上述公式表明它含于由 $\{(1-\iota\sigma^{Mk})(v),\ k=0,\cdots,N-1\}$ 生成的 $\Z_\ell$-子模之中。

\subsection{连通分支与 1-上闭链的空间}

在本节开头我们已指出，最终关心的驯顺上闭链概形 $Z^1(W^0_t,\hatG)$ 是概形 $X_{H,1}=\Hom(W^0_t,H)$ 的一个直和项，其中 $H=\hatG\rtimes W$，$W$ 是 $W_t$ 的有限商，且 $W_t$ 在 $\hatG$ 上的作用通过它分解。反过来，我们下面将看到 $X_{H,1}$ 分解为若干 1-上闭链空间之和，而其中每一个都是连通的。

为简化讨论，我们将在 $\Z[\frac1p]$ 上工作（但由标准的极限论证，下面的一切在它充分大的有限生成子环上同样成立）。在此设定下，$H^0$ 是分裂的，$\pi_0(H)$ 是常值的。把同态前推，得到态射
\[
X_{H,1}=\Hom(W^0_t,H)\xrightarrow{\ \pi\ }\Hom(W^0_t,\pi_0(H))
\]
从 $X_{H,1}$ 到一个有限离散概形，由此得到分解
\[
X_{H,1}=\bigsqcup_{\bar\varphi}X_{H,1,\bar\varphi},\qquad X_{H,1,\bar\varphi}=\pi^{-1}(\{\bar\varphi\}).
\]
下一命题特别表明，每个直和项都非空。

\begin{prop}{3.10}\label{prop:dat-3-10}
对每个 $\bar\varphi:W^0_t\to\pi_0(H)$，存在元素 $\varphi\in X_{H,1,\bar\varphi}(\Z[\frac1p])$，使 $\varphi(W^0_t)$ 有限，并且正规化 $H^0$ 的一个 Borel 对（Borel pair）。此外，若中心 $Z(H^0)$ 在 $\Z[\frac1p]$ 上光滑，则可选取 $\varphi$ 使 $\varphi(W^0_t)$ 正规化 $H^0$ 的一个固定数据（pinning）。
\end{prop}

\begin{proof}
固定 $H^0$ 的一个固定数据 $\varepsilon=(B,T,(X_\alpha)_{\alpha\in\Delta})$，并考虑 $H$ 中底层 Borel 对的正规化子 $\mathcal T:=N_H(B,T)$。由引理 3.1，它是分裂的广义约化群，其中性分支为 $T$，且 $\pi_0(\mathcal T)=\pi_0(H)$。于是可以把迄今的结果应用于 $X_{\mathcal T,1}:=\Hom(W^0_t,\mathcal T)$ 及其直和项 $X_{\mathcal T,1,\bar\varphi}$，后者是 $X_{H,1,\bar\varphi}$ 的闭子概形。

我们来说明 $X_{\mathcal T,1,\bar\varphi}(\Z[\frac1p])\neq\emptyset$。为此先注意映射 $\mathcal T(\Z[\frac1p])\to\pi_0(H)$ 是满的。事实上，这由 [DHKM24, Lemma 2.1] 以及态射 $\mathcal T\to\pi_0(H)$ 是 $T$-主丛这一事实推出。于是把 $\bar\varphi$ 提升为态射 $\varphi:W^0_t\to\mathcal T(\Z[\frac1p])$ 的障碍落在 $H^2_{\mathrm{Ad}\bar\varphi}(W^0_t,\mathcal T(\Z[\frac1p]))$ 中。这里 $\mathrm{Ad}\bar\varphi$ 表示把 $\bar\varphi$ 与 $\mathcal T$ 在 $T=\mathcal T^0$ 上的共轭作用复合所得的 $W^0_t$ 之作用，它通过 $\pi_0(\mathcal T)=\pi_0(H)$ 分解。现在，正如引理 3.6 的证明中已指出的，有
\[
H^2\bigl(W^0_t,\mathcal T(\Z[\frac1p])\bigr)=\bigl[\bigl(\mathcal T(\Z[\frac1p])\bigr)^{\mathrm{Ad}\bar\varphi(s)}\bigr]^{q-1\,\mathrm{Ad}\bar\varphi(\mathrm{Fr})}.
\]
但这个群为零，因为 $\mathrm{Ad}\bar\varphi(\mathrm{Fr})$ 是有限阶自同构，而 $\mathcal T(\Z[\frac1p])$ 是 $p'$-可除的。故 $X_{\mathcal T,1,\bar\varphi}(\Z[\frac1p])\neq\emptyset$。

至此命题的第一个论断已几乎证完，只差像的有限性。然而注意，对任意 $\varphi\in X_{\mathcal T,1,\bar\varphi}(\Z[\frac1p])$，元素 $\sigma=\varphi(s)$ 的阶与 $p$ 互素：它的某个与 $p$ 互素的幂在 $\mathcal T(\overline{\Q})$ 中幂幺，从而等于 $1$。于是，照常用 $F:=\varphi(\mathrm{Fr})$ 记，只需说明存在元素 $c\in\mathcal T(\Z[\frac1p])^\sigma$ 使 $Fc$ 的阶有限，
\ysrc{17}
此时对 $(Fc,\sigma)$ 便定义了所要的同态。为此注意，存在 $n$ 使 $F^n$ 属于 $(T^{\mathrm{Ad}\bar\varphi})^0(\Z[\frac1p])$。但后一个阿贝尔群可除，故存在 $c\in(T^{\mathrm{Ad}\bar\varphi})^0(\Z[\frac1p])$ 使 $c^n=F^n$，于是 $(Fc)^n=1$，如所愿。

至于命题的第二个论断，我们完全照上面进行，只把 $\mathcal T$ 换成固定数据的正规化子 $\mathcal Z$。此时 $H^2_{\mathrm{Ad}\bar\varphi}(W^0_t,\mathcal Z(\Z[\frac1p]))$ 为零，因为 $\mathcal Z^0(\Z[\frac1p])$ 可除，而 $\pi_0(\mathcal Z)$ 是 $p$-群（由我们的光滑性假设），从而是 $p'$-可除的。
\end{proof}

\begin{cor}{3.11}\label{cor:dat-3-11}
设 $\bar\varphi:W^0_t\to\pi_0(H)$，并固定如上一命题中的元素 $\varphi$。以 $\mathrm{Ad}_\varphi$ 记通过 $H$ 内部的共轭得到的 $W_t$ 在 $H^0$ 上的作用。则得到同构
\[
Z^1_{\mathrm{Ad}_\varphi}(W^0_t,H^0)\isom X_{H,1,\bar\varphi},\qquad \eta\mapsto\eta\cdot\varphi.
\]
\end{cor}

这表明 $X_{H,1,\bar\varphi}$ 是 $W^0_t$ 在 $H^0$ 上一个良好的、有限且保持 Borel 对的作用的 1-上闭链空间。在中心光滑的情形，这表明 $X_{H,1,\bar\varphi}$ 是「那个」$F$ 上拟分裂群的驯顺 Langlands 参数空间，该群的 Langlands 对偶为 $H^0\rtimes_{\mathrm{Ad}\varphi}W_F$。

在 [DHKM20] 以及本讲义的第一版中，下述结果是在 $H^0$ 的中心于 $\Z[\frac1p]$ 上光滑的假设下证明的，并在一般情形下作为猜想提出。在本讲义的审稿过程中，一般情形由 Sean Cotner 解决。

\begin{thm}{3.12}（[Cot24]）\label{thm:dat-3-12}
每个非空直和项 $X_{H,1,\bar\varphi}$ 都是连通的。
\end{thm}

Cotner 的途径依赖于上一推论，[DHKM20] 的原始途径亦然。他关于上闭链空间连通性的证明又接近并改进了 5.4 节中给出的下述特殊情形（其中要求作用固定一个固定数据，而不只是保持一个 Borel 对）的证明（特别见注 5.14）。

\begin{thm}{3.13}\label{thm:dat-3-13}
设 $G$ 是驯顺分歧的约化 $p$-进群。则驯顺 Langlands 参数的 $\Z[\frac1p]$-概形 $Z^1(W^0_t,\hatG)$ 连通。
\end{thm}

在 $H^0$ 的中心光滑的附加假设下，借助命题 3.10 的第二部分，这个特殊情形足以得到定理 3.12。

\begin{exam}{3.14}\label{ex:dat-3-14}
设 $G=T$ 是在某个驯顺分歧扩张上分裂的环面。则 $Z^1(W^0_t,\hat T)$ 等同于
\[
\hat T\times\hat T\longrightarrow\hat T,\qquad (F,\sigma)\mapsto F\cdot\mathrm{Fr}\sigma\cdot(s^qF)^{-1}\cdot\bigl(\sigma\cdot{}^s\sigma\cdots{}^{s^{q-1}}\sigma\bigr)^{-1}
\]
的核，因此它具有 $\Z[\frac1p]$ 上可对角化群概形的结构。这样的概形在 $\Z[\frac1p]$ 上连通，当且仅当其特征群的挠部分的阶与 $p$ 互素。若 $s$ 在 $\hat T$ 上的作用平凡（即 $T$ 非分歧），则 $Z^1(W^0_t,\hat T)\simeq\hat T\times\hat T^{\mathrm{Fr}=(\cdot)^q}$，其特征群的挠子群为 $X^*(\hat T)/(q-\mathrm{Fr})X^*(\hat T)$，其阶与 $p$ 互素，如所愿。另一方面，若 $L_s:t\mapsto t(st)^{-1}$ 是 $\hat T$ 的同源，则映射 $Z^1(W^0_t,\hat T)\to\hat T$，$\varphi\mapsto\varphi(s)$ 是满的，其核同构于同源 $L_{s^q}:t\mapsto t(s^qt)^{-1}$ 的核，而后者的次数又与 $p$ 互素。
\ysrc{18}
至于一般情形，考虑次数与 $p$ 互素的同源 $\hat T\twoheadrightarrow\hat T_s\times(\hat T/\hat T^s)$。它诱导态射 $Z^1(W^0_t,\hat T)\to Z^1(W^0_t,\hat T_s)\times Z^1(W^0_t,\hat T/\hat T^s)$，其核与余核都是阶与 $p$ 互素的有限群。于是一般情形由上面两个特殊情形推出。
\end{exam}

\section{归约到驯顺情形与连通分支}

我们现在转向 $Z^1(W^0_F/P,\hatG)$ 的研究，其中 $P$ 是 $P_F$ 的开子群，在 $W_F$ 中可区分（distinguished），且在 $\hatG$ 上平凡作用。

同样，我们研究一个略为一般的设定：固定 $\Lambda$ 上的广义约化群概形 $H$，并考虑仿射概形
\[
X_{H,1P}:=\Hom(W^0_F/P,H),
\]
它表示 $R$-代数上的函子 $R\mapsto\Hom(W^0_F/P,H(R))$。由 Iwasawa 定理 [Iwa55]，存在（非典范的）半直积分解 $W^0_F/P\simeq P_F/P\rtimes W^0_t$。因此我们先研究表示函子 $R\mapsto\Hom(P_F/P,H(R))$ 的仿射概形
\[
Y_{H,1P}:=\Hom(P_F/P,H),
\]
然后使用限制映射 $X_{H,1P}\to Y_{H,1P}$。

\subsection{有限群到 $H$ 的同态}

本段中我们看得更一般些：考虑概形 $Y:=\Hom(\Gamma,H)$，其中 $H$ 是诺特仿射基 $S=\Spec R$ 上的群概形，而 $\Gamma$ 是有限群。共轭诱导群概形 $H$ 在 $Y$ 上的、相对于 $S$ 的作用，它由态射
\[
H\times_S Y\longrightarrow Y,\qquad (h,\phi)\mapsto \mathrm{Ad}_h\circ\phi
\]
给出，其中 $\mathrm{Ad}$ 表示共轭。与该作用相联系的「泛轨道态射」是
\[
\alpha:H\times_S Y\longrightarrow Y\times_S Y,\qquad (h,\phi)\mapsto(\mathrm{Ad}_h\circ\phi,\phi).
\]
这个态射控制通常的轨道态射、中心化子与传输子（transporter）。具体地说：

\begin{itemize}[leftmargin=1.6em,itemsep=0.2em]
\item 对 $\phi\in\Hom(\Gamma,H(R'))$，相应的轨道态射 $H_{R'}\to Y_{R'}$ 是 $Y$-态射 $\alpha$ 沿相应的 $R'$-点 $\Spec R'\to Y$ 的基变换；
\item 对 $\phi,\phi'\in\Hom(\Gamma,H(R'))$，从 $\phi$ 到 $\phi'$ 的传输子是通过沿相应的映射 $S'\to Y\times_S Y$ 拉回 $\alpha$ 而得的 $S':=\Spec(R')$-概形。
\end{itemize}

\begin{lem}{4.1}\label{lem:dat-4-1}
设 $H$ 光滑，且 $\Gamma$ 的阶在 $R$ 中可逆。则 $\Hom(\Gamma,H)$ 光滑，并且一切轨道态射与传输子都光滑。
\end{lem}

\begin{proof}
细节见 [DHKM20, A.1]。由于所有概形在 $S$ 上都有限型，只需验证相应的无穷小提升性质。于是设 $R'$ 是 $R$-代数，$I\subset R'$ 是平方零理想。为证明 $Y$ 在 $S$ 上光滑，需说明任一 $\phi_0:\Gamma\to H(R'/I)$ 都能提升为 $\phi:\Gamma\to H(R')$。但显式计算表明，此事的障碍落在 $H^2_{\mathrm{Ad}\phi_0}(\Gamma,\mathrm{Lie}(H)\otimes_R I)$ 中，而它为零，因为 $\Gamma$ 的阶在 $R$ 中可逆。另一方面，为证明泛轨道态射的光滑性，需说明对任意一对同态 $\phi,\phi':\Gamma\to H(R')$，只要它们的约化
\ysrc{19}
$\phi_0,\phi'_0:\Gamma\to H(R'/I)$ 在某个 $h_0\in H(R'/I)$ 下共轭，就存在 $h_0$ 的提升 $h\in H(R')$，它把 $\phi$ 共轭为 $\phi'$。同样，显式计算表明障碍落在 $H^1_{\mathrm{Ad}\phi_0}(\Gamma,\mathrm{Lie}(H)\otimes_R I)$ 中，而它为零。
\end{proof}

由于轨道态射光滑，它们特别是开的。在 $R$ 为代数闭域的情形，可取 $H(R)$-轨道代表元的集合 $\Phi\subset Y(R)$。于是得到分解 $Y=\coprod_{\phi\in\Phi}H\cdot\phi$，其中 $H\cdot\phi$ 是过 $\phi$ 的轨道态射的像。特别地，$\Phi$ 必须有限。此外，轨道态射诱导同构 $H/C_H(\phi)\isom H\cdot\phi$，其中 $C_H(\phi)$ 是中心化子。

据我们所知，文献中没有关于一般基环 $R$ 的这类一般性结果。不过，下面这个结果足以满足我们的需要。

\begin{thm}{4.2}\label{thm:dat-4-2}
设 $H$ 是广义约化群，$R$ 是数域的正规子环，且 $\Gamma$ 的阶在 $R$ 中可逆。
\begin{enumerate}[label=(\arabic*),leftmargin=2.4em,itemsep=0.25em]
\item 范畴商 $Y\sslash H^0=\Spec(\OO_Y)^{H^0}$ 在 $R$ 上有限 étale，并代表 étale 拓扑下的商层 $Y/H^0$。
\item 对任意 $\phi\in Y(R')$，étale 拓扑下的商层 $H^0_{R'}/C_{H^0}(\phi)$ 可由一个 $R'$-概形表示，而轨道态射 $H^0_{R'}\to Y_{R'}$ 把它等同于 $Y_{R'}$ 的一个开且闭的子概形。
\item 存在有限正规扩张 $R'/R$ 与有限集 $\Phi\subset Y(R')$，使轨道映射诱导同构 $\coprod_{\phi\in\Phi}H^0_{R'}/C_{H^0}(\phi)\isom Y_{R'}$。
\end{enumerate}
\end{thm}

\begin{proof}
对 (1)，我们引用 [DHKM20, Thm. A.7]，只指出两个主要成分。关键点是证明：当 $R$ 是严格 Hensel 离散赋值环时，从 $Y(R)/H^0(R)$ 所联系的常值层到商层 $Y/H^0$ 的典范映射是同构。换言之，对任意严格 Hensel 局部 $R$-代数 $R'$，映射 $Y(R)/H^0(R)\to Y(R')/H^0(R')$ 应当是双射。若映射 $R\to R'$ 是局部的，则有分解
\[
Y(R)/H^0(R)\isom Y(k_R)/H^0(k_R)\longrightarrow Y(k_{R'})/H^0(k_{R'})\simeq Y(R')/H^0(R'),
\]
其中同构来自传输子的光滑性与 Artin 逼近。于是只剩域扩张的情形，这可以用 V. Lafforgue 的伪特征标处理。否则我们有分解 $Y(R)/H^0(R)\to Y(\overline K)/H^0(\overline K)\to Y(k_{R'})/H^0(k_{R'})\simeq Y(R')/H^0(R')$，其中 $\overline K$ 是 $R$ 的分式域的代数闭包，并用 Bruhat–Tits 理论得到第一个映射的双射性。正是在这一部分我们用到了 $R$ 是一维正则环。

对 (2)，只需对与 $Y$ 的恒等相对应的泛同态 $\phi_{\mathrm{univ}}$ 证明。其轨道态射为 $\alpha:H^0\times_S Y\to Y\times_S Y$。令 $C_{H^0}(\phi_{\mathrm{univ}})\to Y$ 为 $Y$ 上的泛中心化子，即 $\alpha$ 沿对角线的拉回。则 $\alpha$ 诱导同构
\[
(H^0\times_S Y)/C_{H^0}(\phi_{\mathrm{univ}})\isom Y\times_{Y/H^0}Y,
\]
其中所有商都是 étale 拓扑下的层论商。由 (1) 知目标可由仿射概形表示，故源也可表示。此外，态射 $Y\times_{Y/H^0}Y\to Y\times_S Y$ 是
\[
Y/H^0=Y/H^0\times_{Y/H^0}Y/H^0\longrightarrow Y/H^0\times_S Y/H^0
\]
的基变换，而后者是开闭浸入，因为 $Y/H^0$ 在 $S$ 上有限 étale。

对 (3)，我们引用 [DHKM20, Thm. A.9]，只说明问题与解决它的主要工具。可以先用引理 3.2
\ysrc{20}
把 $R$ 换成有限扩张，使 $H$ 在 $R$ 上分裂。再利用 (1)，把 $R$ 换成为使 $\OO(Y)^{H^0}$ 分裂的有限 étale 扩张。于是问题在于找到 $R'$，使 $(Y/H^0)(R)$ 的任一点 $[\phi]$ 都能提升为某个代表元 $\phi\in Y(R')$。由 $Y\to Y/H^0$ 的光滑性，我们当然可以找到忠实 étale 于 $R$ 的 $R'$，使该性质成立。因此，在把 $R$ 换成有限扩张之后，可以找到 $R'$ 使 $\Spec R'\to\Spec R$ 是 Zariski 覆盖，即 $R'=\prod_{i=1}^rR_i$，其中每个 $R_i$ 都是 $R$ 的局部化。于是问题就在于从该 Zariski 覆盖过渡到某个有限平坦覆盖。在 loc.cit. 中，这是借助强逼近定理完成的。
\end{proof}

\begin{rem}{4.3}
若假设 $H$ 是约化群，则它在 $R$ 上某个忠实 étale 的 $R'$ 上分裂，而分裂的 $H_{R'}$ 按分类理论可下降到 $\Spec\Z$ 上。此时可用该定理证明：只要环 $R$ 在 $\Z$ 上平坦或者含有一个域，其结论对 $R$ 依然成立。
\end{rem}

\subsection{限制到 $P_F$}

现在我们重新拾起对 $X_{H,1P}=\Hom(W^0_F/P,H)$ 的研究，并回到本节开头的记号。记住基环 $\Lambda$ 是 $\Z[\frac1p]$ 的有限整扩张，因此上面的定理 4.2 适用于 $Y_{H,1P}=\Hom(P_F/P,H)$。在必要时扩大 $\Lambda$，我们可以假设存在有限集 $\Phi\subset Y_{H,1P}(\Lambda)$，使 $Y_{H,1P}$ 分解为轨道之和
\[
Y_{H,1P}=\coprod_{\varphi\in\Phi}Y_{H,[\varphi]},
\qquad\text{其中}\quad
H^0/C_{H^0}(\varphi)\isom H^0\cdot\varphi=:Y_{H,[\varphi]}.
\]
沿限制态射 $X_{H,1P}\to Y_{H,1P}$ 拉回，便得到分解
\[
X_{H,1P}=\coprod_{\varphi\in\Phi}X_{H,[\varphi]},
\qquad\text{其中}\quad
H^0\times^{C_{H^0}(\varphi)}X_{H,\varphi}\isom X_{H,[\varphi]},
\]
而
\[
X_{H,\varphi}=\Hom(W^0_F,H)_\varphi:=\{\phi\in X_{H,1P},\ \phi|_{P_F}=\varphi\}.
\]
上面的收缩积表示 $H^0\times X_{H,\varphi}$ 对 $C_{H^0}(\varphi)$ 的对角作用的商。这个商先作为 étale 拓扑下的层来定义，而上面的同构（在层论框架下相当形式）表明它确实可表示。

当然，$X_{H,\varphi}$ 可能是空的。我们说 $\varphi$ 是容许的（admissible），如果 $X_{H,\varphi}\neq\emptyset$。

对 $\varphi\in\Phi$，记 $H_\varphi:=N_H(\varphi(P_F))$ 为 $\varphi$ 在 $H(\Lambda)$ 中的像在 $H$ 中的概形论正规化子。由引理 3.2，它仍是 $\Lambda$ 上的广义约化群概形。注意 $H_\varphi(\Lambda)$ 含 $\varphi(P_F)$，于是如上我们有一个 $\Lambda$-概形 $X_{H_\varphi,\varphi}=\Hom(W^0_F/P,H_\varphi)_\varphi$。此外，由于 $P_F$ 在 $W_F$ 中可区分，任意 $\phi:W^0_F/P\to H(R)$ 都经过 $H_\varphi(R)$，从而闭浸入 $H_\varphi\into H$ 诱导同构
\[
X_{H_\varphi,\varphi}=\Hom(W^0_F/P,H_\varphi)_\varphi
\isom\Hom(W^0_F/P,H)_\varphi=X_{H,\varphi}.
\]
现在如 [Iwa55] 那样固定一个分裂 $W^0_F/P=W^0_t\ltimes P_F/P$。这给出限制映射
\[
X_{H_\varphi,\varphi}=\Hom(W^0_F/P,H_\varphi)_\varphi
\longrightarrow\Hom(W^0_t,H_\varphi)=X_{H_\varphi,1}.
\]

\begin{lem}{4.4}\label{lem:dat-4-4}
上述限制映射是开闭嵌入。
\end{lem}
\ysrc{21}
\begin{proof}
考虑态射 $X_{H_\varphi,1}\to Y_{H_\varphi,1P}\times_\Lambda Y_{H_\varphi,1P}$，它把 $(F,\sigma)$ 映为 $(F^\varphi,\sigma^\varphi)$，其中 $F^\varphi=\mathrm{Ad}_F\circ\varphi\circ\mathrm{Ad}_{\mathrm{Fr}}^{-1}$，$\sigma^\varphi=\mathrm{Ad}_\sigma\circ\varphi\circ\mathrm{Ad}_s^{-1}$。于是 $X_{H_\varphi,\varphi}$ 等同于该态射沿对角嵌入 $\Spec\Lambda\xrightarrow{(\varphi,\varphi)}Y_{H_\varphi,1P}\times_\Lambda Y_{H_\varphi,1P}$ 的拉回，故只需说明这个对角嵌入是开闭浸入。但映射 $\Spec\Lambda\xrightarrow{\varphi}Y_{H_\varphi,1P}$ 诱导同构 $\Spec\Lambda\isom Y_{H_\varphi,[\varphi]}$，因为 $(H_\varphi)^0=C_H(\varphi)^0$。由于 $Y_{H_\varphi,[\varphi]}$ 在 $Y_{H_\varphi,1P}$ 中开闭，这说明 $\varphi$ 是开闭浸入，从而 $(\varphi,\varphi)$ 也是。
\end{proof}

\subsection{定理 2.1 与 2.2 的证明}

这两个定理都是下面这个定理的推论。

\begin{thm}{4.5}\label{thm:dat-4-5}
$\Lambda$-概形 $X_{H,1P}=\Hom(W^0_F/P,H)$ 是约化的、平坦的、相对维数为 $\dim H$ 的局部完全交。对任意素数 $\ell\neq p$，其函数环 $R_{H,1P}$ 是 $\ell$-进分离的，且泛同态 $\varphi_{H,1P}:W^0_F/P\to R_{H,1P}$ 唯一地延拓为泛的 $\ell$-进与 ind-$\ell$-进连续同态 ${}^\ell\varphi_{H,1P}:W_F/P\to H(R_{H,1P}\otimes\Z_\ell)$。
\end{thm}

\begin{proof}
对容许的 $\varphi\in\Phi$，$\Lambda$-概形 $X_{H,\varphi}$ 是 $X_{H_\varphi,1}$ 的直和项，故由命题 3.3 它是平坦的、相对维数为 $\dim H_\varphi$ 的局部完全交，而其函数环 $R_{H,\varphi}$ 由引理 3.8 是 $\ell$-进分离的。泛同态 $\varphi_{H,\varphi}:W^0_F/P_F\to H(R_{H,\varphi})$ 相对于为限制映射 $X_{H_\varphi,\varphi}\to X_{H_\varphi,1}$ 所选取的分裂 $W_F/P=W_t\ltimes(P_F/P)$ 具有形式 $\varphi_{H_\varphi,1}\ltimes\varphi$。因此，用命题 3.9 的记号，我们看到 ${}^\ell\varphi_{H,\varphi}:={}^\ell\varphi_{H_\varphi,1}\ltimes\varphi$ 是 $\varphi_{H,\varphi}$ 的泛连续延拓。

现在，由同构 $X_{H,[\varphi]}\simeq H^0\times^{C_{H^0}(\varphi)}X_{H,\varphi}$ 可知，$\Lambda$-概形 $X_{H,[\varphi]}$ 是平坦的、相对维数为 $\dim H$ 的局部完全交，且其函数环 $R_{H,[\varphi]}=(\OO_{H^0}\otimes_\Lambda R_{H,\varphi})^{C_{H^0}(\varphi)}$ 是 $\ell$-进分离的。此外，泛同态 $\varphi_{H,[\varphi]}$ 由如下公式给出：
\[
\varphi_{H,[\varphi]}(w)\colon\ \OO_H\longrightarrow\OO_{H^0}\otimes_\Lambda\OO_H
\xrightarrow{\ \mathrm{Id}\otimes\varphi_{H,\varphi}(w)\ }\OO_{H^0}\otimes_\Lambda R_{H,\varphi},
\qquad \forall w\in W_F,
\]
其中第一个映射是 $(g,h)\mapsto ghg^{-1}$ 的对偶。把一切基变换到 $\Z_\ell$ 并把 $\varphi_{H,\varphi}$ 换成 ${}^\ell\varphi_{H,\varphi}$，就得到所要的 ${}^\ell\varphi_{H,[\varphi]}$。最后只需利用分解 $X_{H,1P}=\coprod_{\varphi\in\Phi}X_{H,[\varphi]}$ 来完成定理的证明。
\end{proof}

我们要指出，这里给出的证明比 [DHKM20] 的证明明显简单。简化来自系统地使用 Iwasawa 分裂（如上一个引理中那样），这借自 [Zhu25] 第 3.1 节。

\subsection{$\overline{\Z}[\frac1p]$ 上的连通分支与驯顺 1-上闭链空间}

我们保留上面同样的 $\Phi$ 与 $X_{H,\varphi}=\Hom(W^0_F/P,H)_\varphi$ 的记号，并为了方便把一切基变换到 $\overline\Z[\frac1p]$。与 3.5 小节一样，我们可以用态射
\[
X_{H,\varphi}=\Hom(W^0_F,H)_\varphi=\Hom(W^0_F,H_\varphi)_\varphi
\longrightarrow\Hom(W^0_F,\pi_0(H_\varphi))_\varphi
\]
得到更细的分解
\[
X_{H,\varphi}=\coprod_{\bar\phi}X_{H,\varphi,\bar\phi}
\]
\ysrc{22}
其中 $\bar\phi$ 取遍有限集 $\Hom(W^0_F,\pi_0(H_\varphi))_\varphi$。同样，我们说对 $(\varphi,\bar\phi)$ 是容许的，如果 $X_{H,\varphi,\bar\phi}\neq\emptyset$。

\begin{thm}{4.6}\label{thm:dat-4-6}
对每个容许对 $(\varphi,\bar\phi)$，概形 $X_{H,\varphi,\bar\phi}$ 连通。此外，存在 $\phi\in X_{H,\varphi,\bar\phi}(\Z[\frac1p])$，使 $\phi(W^0_F)$ 有限，且正规化 $(H_\varphi)^0=C_H(\varphi)^0$ 的一个 Borel 对。若 $(H_\varphi)^0$ 的中心在 $\Z[\frac1p]$ 上光滑，则甚至可选 $\phi$ 使其保持 $(H_\varphi)^0$ 的一个固定数据（pinning）。
\end{thm}

注意，对任意 $(\varphi,\bar\phi)$ 以及如命题中那样的 $\phi$，我们得到同构
\begin{equation*}\tag*{(4.1)}
Z^1_{\mathrm{Ad}\phi}\bigl(W^0_t,(H_\varphi)^0\bigr)
\xrightarrow{\ \sim\ }X_{H,\varphi,\bar\phi},\qquad \eta\mapsto\eta\cdot\phi,
\end{equation*}
我们可以把它看成「归约到驯顺参数」。

\begin{proof}
由引理 4.4 我们知道，在取定分裂 $W^0_F/P_F=W^0_t\ltimes(P_F/P)$ 之后，限制映射 $X_{H,\varphi,\bar\phi}\to X_{H_\varphi,1,\bar\phi}$ 是开闭嵌入。由 Cotner 的定理 3.12，$X_{H_\varphi,1,\bar\phi}$ 连通。由于 $X_{H,\varphi,\bar\phi}\neq\emptyset$，这个限制映射是同构，从而 $X_{H,\varphi,\bar\phi}$ 连通。由于 $\varphi(P_F)$ 中心化 $(H_\varphi)^0$，其余论断由命题 3.10 推出。
\end{proof}

\begin{rem}{4.7}
上述定理中 $\phi$ 的存在性就是 [DHKM20] 的定理 3.4，其证明可以说比上面五行论证复杂得多。这里的简化之一来自引理 4.4，而它在 [DHKM20] 中被忽略了。另一个「简化」是我们在这里用了 Cotner 的连通性结果。注意 [Cot24] 的主要结果依赖归约到驯顺参数，但使用驯顺版本（即 [Cot24] 的定理 4.7），我们可以避免任何循环论证。

下面给出 $\phi$ 存在性的一个更自足的证明，它避开了 Cotner 的连通性定理（其证明长达 8 页），但仍简化了原来的证明。
\end{rem}

\begin{proof}[归约到驯顺参数的另一个证明]
沿用上一证明的记号，我们先走命题 3.10 证明的最初几步：固定 $(H_\varphi)^0$ 的一个固定数据 $\varepsilon=(B,T,(X_\alpha)_{\alpha\in\Delta})$，并考虑 $(H_\varphi)^0$ 中底层 Borel 对的正规化子 $\mathcal T:=N_{H_\varphi}(B,T)$。注意 $\varphi(P_F)$ 含于 $\mathcal T(\Lambda)$ 且在其中正规，故 $\mathcal T=\mathcal T_\varphi$。此外，$\pi_0(\mathcal T)=\pi_0(H_\varphi)$。因此限制映射 $X_{\mathcal T,\varphi,\bar\phi}\to X_{\mathcal T,1,\bar\phi}$ 是开闭嵌入，而由例 3.14 知 $X_{\mathcal T,1,\bar\phi}$ 连通。于是只需证明概形 $X_{\mathcal T,\varphi,\bar\phi}$ 非空。为此，取代数闭域 $L$ 使 $X_{H,\varphi,\bar\phi}(L)\neq\emptyset$，并取 $\phi_0\in X_{H,\varphi,\bar\phi}(L)$。对任意 $w\in W^0_F$，存在 $h_w\in(H_\varphi)^0(L)$ 使 $\mathrm{Ad}_{h_w}\circ\mathrm{Ad}_{\phi_0(w)}$ 属于 $\mathcal T(L)$。这些 $h_w$ 只在乘以 $\mathcal T(L)$ 的意义下确定，而使 $w\mapsto h_w\phi_0(w)$ 成为同态的障碍落在 $H^2_{\mathrm{Ad}\bar\phi}(W^0_t,\mathcal T(L))$ 中（细节见 [DHKM20, Prop. 3.7]）。如命题 3.10 的证明中那样，这个 $H^2$ 为零，因为 $\mathcal T(L)$ 是可除群。在 $Z((H_\varphi)^0)$ 光滑的情形，改用固定数据的正规化子 $\mathcal Z$ 代替 $\mathcal T$ 作类似论证。
\end{proof}

\begin{rem}{4.8}
在 $Z(H^0)$ 于 $\Z[\frac1p]$ 上光滑的假设下，[DHKM20] 证明了 $X_{H,\varphi,\bar\phi}$ 的连通性。确实，loc. cit. 的引理 3.11 表明，此时 $Z((H_\varphi)^0)$ 也在 $\Z[\frac1p]$ 上光滑。于是，如上面的替代证明中那样，存在 $\phi$ 使 $\mathrm{Ad}_\phi$ 保持 $(H_\varphi)^0$ 的一个固定数据，而 $X_{H,\varphi,\bar\phi}$ 的连通性随即由定理 3.13 推出。
\end{rem}
\ysrc{23}
\begin{rem}{4.9}
虽然 Cotner 能够绕过保持固定数据的上闭链的存在性，但当我们从形如 $H={}^LG$ 的 $H$ 出发、并且坚持 $\bar\phi$ 与典范映射 $W_F\to\pi_0({}^LG)$ 相容时，是否总能归约到这种情形，仍是一个未解决的问题。
\end{rem}

为了参数化 $X_{H,1P}$ 的连通分支，必须把中心化子 $C_{H^0}(\varphi)$ 可能的非连通性考虑在内。这个群概形通过共轭作用在集合 $\pi_\varphi:=\Hom(W^0_F,\pi_0(H_\varphi))_\varphi$ 上。记 $C_{H^0}(\varphi)_{\bar\phi}$ 为元素 $\bar\phi$ 的稳定化子，记 $\bar\pi_\varphi$ 为 $\pi_0(C_{H^0}(\varphi))$ 在 $\Hom(W^0_F,\pi_0(H_\varphi))_\varphi$ 上的轨道的一组代表元。于是我们得到分解
\begin{equation*}\tag*{(4.2)}
X_{H,1P}=\coprod_{\varphi\in\Phi}\ \coprod_{\bar\phi\in\bar\pi_\varphi}
H^0\times^{C_{H^0}(\varphi)_{\bar\phi}}X_{H,\varphi,\bar\phi}.
\end{equation*}
由上述结果，这给出 $X_{H,1P}$ 分解为连通分支。

\subsection{表示论方面预期的镜像性质}

设 $G$ 是 $F$ 上的拟分裂约化群，其 $L$-群为 ${}^LG=\hatG\rtimes W_F$。在 $\overline\Z[\frac1p]$ 上工作，我们得到了分解
\begin{equation*}\tag*{(4.3)}
Z^1(W^0_F,\hatG):=\varinjlim_{P}Z^1(W^0_F/P,\hatG)
=\coprod_{\varphi,\phi}\hatG\times^{C_{\hatG}(\varphi)_{\bar\phi}}
Z^1_{\mathrm{Ad}\phi}\bigl(W^0_t,C_{\hatG}(\varphi)^0\bigr),
\end{equation*}
其中
\begin{itemize}[leftmargin=1.6em,itemsep=0.2em]
\item $\varphi$ 取遍 $Z^1(P_F,\hatG(\overline\Z[\frac1p]))$ 中一组代表元所成的集合 $\Phi$，这些代表元代表 $\hatG$-轨道的「野惯性参数」；
\item 给定 $\varphi$ 之后，$\phi$ 取遍一个有限的上闭链集合 $W_F\to\hatG(\overline\Z[\frac1p])$，这些上闭链延拓 $\varphi$ 并正规化 $C_{\hatG}(\varphi)^0$ 的一个 Borel 对。
\end{itemize}
回忆我们说对 $(\varphi,\phi)$ 是「容许的」，如果相应的直和项非零；此时该直和项现已知道是连通的。在 $\hatG$ 于 $\overline\Z[\frac1p]$ 上有光滑中心的情形（例如 $G_{\mathrm{der}}$ 单连通），可以选取 $\phi$ 使其正规化 $C_{\hatG}(\varphi)^0$ 的一个固定数据。设 $G_{\varphi,\phi}$ 是拟分裂的驯顺分歧群，其对偶群 $\widehat{G_{\varphi,\phi}}$ 就是带有作用 $\mathrm{Ad}_\phi$ 的 $C_{\hatG}(\varphi)^0$。于是 $\phi$ 诱导嵌入
\[
\zeta_{\varphi,\phi}:{}^LG_{\varphi,\phi}=C_{\hatG}(\varphi)^0\rtimes W_F
\longrightarrow{}^LG=\hatG\rtimes W_F,
\]
它与到 $W_F$ 的投影相容，即是一个 $L$-群的态射。

这些事实提示了光滑 $\overline\Z[\frac1p]G(F)$-模范畴 $\mathrm{Rep}_{\overline\Z[\frac1p]}(G(F))$ 上如下的镜像性质。

\medskip
\noindent\hangindent=2.2em\hangafter=1
(1)\ 应当存在作为全子范畴之积的分解
\[
\mathrm{Rep}_{\overline\Z[\frac1p]}G(F)=\prod_{\varphi,\phi}\mathrm{Rep}_{\varphi,\phi}G(F),
\]
其特征性质是：$\pi\in\mathrm{Irr}_{\C}G(F)$ 属于 $\mathrm{Rep}_{\varphi,\phi}G(F)$ 当且仅当 Langlands 参数 $\phi_\pi$ 满足 $(\phi_\pi)|_{P_F}\sim\varphi$ 且 $\overline{\phi_\pi}\sim\overline{\phi}$。换言之，由这些表示切出的 Bernstein 中心 $\mathfrak Z_{\C}(G(F))$ 的幂等元 $e_{\varphi,\phi}$ 应当属于 $\mathfrak Z_{\overline\Z[\frac1p]}(G(F))$。此外，每个因子 $\mathrm{Rep}_{\varphi,\phi}G(F)$ 都应当是 $\mathrm{Rep}_{\overline\Z[\frac1p]}(G(F))$ 的一个「稳定块」
\ysrc{24}
，即在环 $\mathfrak Z_{\overline\Z[\frac1p]}(G(F))\cap\mathfrak Z^{st}_{\C}(G(F))$ 中 $e_{\varphi,\phi}$ 应当是本原的（primitive）这个意义下；其中 $\mathfrak Z^{st}_{\C}(G(F))$ 表示 $\mathfrak Z_{\C}(G(F))$ 中这样的元素所成的子环：它们在同一个（设想的）$L$-包中的任意两个不可约表示上作用为同一个标量。

\smallskip
\noindent\hangindent=2.2em\hangafter=1
(2)\ 当 $C_{\hatG}(\varphi)$ 连通时，应当存在忠实嵌入
\[
\zeta^*_{\varphi,\phi}:\ \mathrm{Rep}_1(G_{\varphi,\phi}(F))
\hookrightarrow \mathrm{Rep}_{\varphi,\phi}G(F),
\]
它在每侧取定 Whittaker 数据之后，在复不可约表示上诱导与 $\zeta_{\varphi,\phi}$ 相联系的 Langlands 转移映射，并把 $G_{\varphi,\phi}(F)$ 的 Whittaker 空间的深度 $0$ 因子送到 $G$ 的 Whittaker 空间的 $(\varphi,\phi)$-因子。

在 $C_{\hatG}(\varphi)^0$ 非连通的情形，嵌入 $\zeta^*_{\varphi,\phi}$ 的源应当是 $\mathrm{Rep}_1(G_{\varphi,\phi}(F)')$，其中 $G_{\varphi,\phi}(F)'$ 是 $\pi_0(C_{\hatG}(\varphi)_{\bar\phi})$ 的某个商被 $G_{\varphi,\phi}(F)$ 所作的适当扩张（在略为不同的语境下的讨论见 [Dat17] 第 2.1.3 节与例 2.1.4）。

此外，扩张 LLC 的形式体系提示如下加细。

\smallskip
\noindent\hangindent=2.2em\hangafter=1
(3)\ 记 $S_{\varphi,\phi}$ 为有限阿贝尔群 $\pi_0(Z(C_{\hatG}(\varphi)^0)^{\phi(W_F)}/Z(\hatG)^{W_F})$ 的特征标所成的阿贝尔群。则应当存在进一步的分解
\[
\mathrm{Rep}_{\varphi,\phi}G(F)=\coprod_{\alpha\in S_{\varphi,\phi}}\mathrm{Rep}^\alpha_{\varphi,\phi}G(F),
\]
其特征性质是：$\pi\in\mathrm{Irr}_{\C}G(F)$ 属于 $\mathrm{Rep}^\alpha_{\varphi,\phi}G(F)$ 当且仅当扩张 Langlands 参数 $(\phi_\pi,\varepsilon_\pi)$ 满足 $(\phi_\pi)|_{P_F}\sim\varphi$、$\phi_\pi\sim\phi$，且 $(\varepsilon_\pi)|_{S_{\varphi,\phi}}$ 是 $\alpha$-同型的（这里 $\varepsilon_\pi$ 是群 $S_{\phi_\pi}:=\pi_0(C_{\hatG}(\phi_\pi)/Z(\hatG)^{W_F})$ 的不可约表示，而限制沿包含 $Z(C_{\hatG}(\varphi)^0)^{\phi(W_F)}\subset C_{\hatG}(\phi_\pi)$ 所诱导的映射 $S_{\varphi,\phi}\to S_{\phi_\pi}$ 来取）。换言之，由这些表示切出的 Bernstein 中心 $\mathfrak Z_{\C}(G(F))$ 的幂等元 $e^\alpha_{\varphi,\phi}$ 应当属于 $\mathfrak Z_{\overline\Z[\frac1p]}(G(F))$。此外，每个因子 $\mathrm{Rep}^\alpha_{\varphi,\phi}G(F)$ 都应当是 $\mathrm{Rep}_{\overline\Z[\frac1p]}(G(F))$ 的一个块。

\smallskip
\noindent\hangindent=2.2em\hangafter=1
(4)\ 当 $C_{\hatG}(\varphi)$ 连通时，应当存在范畴等价
\[
\mathrm{Rep}_1(G^\alpha_{\varphi,\phi}(F))\xrightarrow{\ \sim\ }\mathrm{Rep}^\alpha_{\varphi,\phi}G(F),
\]
它诱导不可约表示的通常转移。这里 $\alpha\in S_{\varphi,\phi}$ 经由 Kottwitz 同构被视为映射 $H^1(F,G_{\varphi,\phi})\to H^1(F,G)$ 的核中的元素，该映射由包含 $Z(C_{\hatG}(\varphi)^0)^{\phi(W_F)}\supset Z(\hatG)^{W_F}$ 推出，而 $G^\alpha_{\varphi,\phi}$ 是 $G_{\varphi,\phi}$ 的相应纯内型。

在 $C_{\hatG}(\varphi)$ 非连通的情形，群 $G^\alpha_{\varphi,\phi}(F)$ 应换为该群被 $\pi_0(C_{\hatG}(\varphi)_{\bar\phi})$ 的某个商所作的适当扩张。

在 (2) 与 (4) 中，$\mathrm{Rep}_1$ 对应于参数对 $(1,1)$，并且由于 (1)，它应当是 $G^\alpha_{\varphi,\phi}(F)$ 的深度 $0$ 范畴。换言之，这些预测是关于一个函子性过程的：把 $G(F)$ 的表示论归约到辅助的驯顺分歧群的深度 $0$ 表示论。
\ysrc{25}
\begin{rem}{4.10}
即使 $\hatG$ 没有光滑中心，我们也预期该分解存在，范畴等价也存在，尽管可能涉及一些扭（twist）。关于连通性，注意当 $p>2$ 时，对经典群 $C_{\hatG}(\varphi)$ 总是连通的。
\end{rem}

以下是我们目前所知道的。预测的分解 (1) 现在可由 Scholze（2025）关于 Fargues–Scholze 谱作用的 motivic 版本的近期工作推出。在此之前，(1) 与 (3) 中的分解已在 [Dat18] 中于 $G$ 为「非常驯顺」的假设下构造出来，即 $G$ 的一切极大环面都驯顺分歧；此时 $C_{\hatG}(\varphi)$ 连通，而且实际上是 $\hatG$ 的 Levi 子群。那里的构造深受 Kaletha 对这类群的超尖点包的构造 [Kal19] 的启发。特别地，它与他提出的 Langlands 对应相容。因此主要工具是 Yu 的泛特征标理论 [Yu01] 与 Fintzen 关于 Yu 构造的穷尽性结果 [Fin21]。(4) 中那样的范畴等价之构造出现在作者与 Fintzen 的进行中工作里。深度 $0$ 范畴不可分解这一事实在 [DL22] 中对驯顺分歧群得到证明。对 $\mathrm{GL}_n$，(1) 那样的分解由 Bushnell–Kutzko 型理论推出，而且不需要 (3) 那样的加细。此时范畴等价可由 Chinello 的工作 [Chi18] 推出。不过，这些等价与 LLC 的相容性远未清楚。

\begin{rem}{4.11}
这些预测显然与 CLLC 哲学相容，只有两点例外：我们使用的系数 $\Z[\frac1p]$，以及我们坚持在阿贝尔范畴中工作。更多讨论见 4.7 小节。
\end{rem}

\subsection{$\Z_\ell$ 上的进一步分解}

现在固定素数 $\ell\neq p$，为简单起见在 $\Z_\ell$ 上工作。同样，只要固定深度 $P\subset P_F$，一切都在 $\Z_\ell$ 的某个足够大的扩张上成立。记 $I^\ell_F$ 为映射 $I_F\xrightarrow{t_\ell}\Z_\ell(1)$ 的核。它是 $I_F$ 的极大 pro-$\ell'$ 子群。它将扮演迄今由 $P_F$ 扮演的角色。

我们已经看到，$Z^1(W^0_F/P,\hatG)_{\Z_\ell}$ 上的泛 1-上闭链有唯一的 $\ell$-进连续延拓 ${}^\ell\varphi_{\hatG,1P}$ 到 $W_F$。我们断言，${}^\ell\varphi_{\hatG,1P}$ 在 $I^\ell_F$ 上的限制通过 $I^\ell_F$ 的一个依赖于 $P$ 的有限商分解。如果我们用 condensed 参数来定义 $Z^1$，这一点将是显然的。在我们的设定中，可以归约到驯顺情形，在每个几何一般点处取值，并利用 ${}^\ell\varphi_{\hatG,1}(s)$ 在这样的点处的半单部分具有有限阶这一事实，如命题 3.3 的证明中那样。无论如何，通过限制，我们于是得到态射
\[
Z^1(W^0_F,\hatG)_{\Z_\ell}\longrightarrow Z^1(I^\ell_F,\hatG)_{\Z_\ell},
\]
其中右边表示 $I^\ell_F$ 的有限商的上闭链空间的余极限。由于这样的有限商的阶在 $\Z_\ell$ 中可逆，我们可以应用 4.1 小节的结果。于是设 $\Phi_\ell\subset Z^1(I^\ell_F,\hatG(\overline\Z_\ell))$ 是一组代表 $\hatG(\overline\Z_\ell)$-共轭类的代表元。对每个 $\varphi_\ell\in\Phi_\ell$，记 $Z^1(W^0_F,\hatG)_{\Z_\ell,\varphi_\ell}$ 为满足 $({}^\ell\varphi_{\hatG})|_{I^\ell_F}=\varphi_\ell$ 的闭子集。通过拉回，我们得到分解
\[
Z^1(W^0_F,\hatG)_{\Z_\ell}=\coprod_{\varphi_\ell\in\Phi_\ell}
\hatG\times^{C_{\hatG}(\varphi_\ell)}Z^1(W^0_F,\hatG)_{\Z_\ell,\varphi_\ell}.
\]
\ysrc{26}
如上面几个小节那样论证，我们可以把这个分解加细如下。对每个 $\phi_\ell$，存在有限子集 $\Sigma_{\phi_\ell}\subset Z^1(W^0_F,\hatG(\Z_\ell))_{\phi_\ell}$，使每个 $\varphi\in\Sigma_{\phi_\ell}$ 都有有限的像，且正规化 $C_{\hatG}(\phi_\ell)^0$ 的一个 Borel 对，并且有
\begin{equation*}\tag*{(4.4)}
Z^1(W^0_F,\hatG)_{\Z_\ell}
=\coprod_{\phi_\ell\in\Phi_\ell}\ \coprod_{\varphi\in\Sigma_{\phi_\ell}}
\hatG\times^{C_{\hatG}(\phi_\ell)_\varphi}
Z^1_{\mathrm{Ad}\varphi}\bigl(W^0_t,C_{\hatG}(\phi_\ell)^0\bigr)_1,
\end{equation*}
记号与 4.4 小节末相同。这里要注意，$W_F$ 在 $C_{\hatG}(\phi_\ell)^0$ 上的作用 $\mathrm{Ad}_\varphi$ 是「驯顺 $\ell$-分歧的」，意思是它驯顺分歧，且惯性通过一个 $\ell$-群商作用。下标 $1$ 表示我们看的是与 $\phi_\ell=1_{I^\ell_F}$ 相联系的闭子集。同样，在中心 $Z(\hatG)$ 于 $\Z_\ell$ 上光滑的情形，我们可以选取 $\Sigma_{\phi_\ell}$ 使每个 $\varphi\in\Sigma_{\phi_\ell}$ 都正规化 $C_{\hatG}(\phi_\ell)^0$ 的一个固定数据。于是，粗略地说，这个分解表明 $Z^1(W^0_F,\hatG)_{\Z_\ell}$ 如何由驯顺 $\ell$-分歧群的驯顺 $\ell$-分歧参数空间构成。下面这个定理将在第 5.5 节证明，它说明上述最后一个分解就是连通分支的分解。

\begin{thm}{4.12}\label{thm:dat-4-12}
设 $\hatG$ 带有保持 Borel 对的有限驯顺 $\ell$-分歧作用。则 $Z^1(W^0_t,\hatG)_{\Z_\ell,1}$ 连通。
\end{thm}

\subsection{$\Z_\ell$ 系数下的镜像性质}

上述模式再次提示，对 $F$ 上的拟分裂约化群 $G$，关于范畴 $\mathrm{Rep}_{\Z_\ell}(G(F))$ 有一个相应的镜像模式，仍沿 4.5 小节中的 (1)–(4) 各条。我们不逐一重复，但要指出 (3) 中的一个重要差别：我们不应期望因子范畴 $\mathrm{Rep}^{\alpha}_{\overline\Z_\ell,(\phi_\ell,\varphi)}(G(F))$ 是不可分解的。这与 Bernstein 中心里「不稳定」本原幂等元的存在有关，意思是这些幂等元在某个 $L$-包中消灭了一个真子集的表示。例如，取 $G=\mathrm{GSp}_4$ 且 $\phi_\ell=1$（从而也有 $\varphi=1$ 与 $\alpha=1$）。对大的 $\ell$，相应的因子 $\mathrm{Rep}_{\overline\Z_\ell,1}(G(F))$ 由幂幺表示组成，并且像特征 $0$ 情形那样，分解为 Iwahori 块与由幂幺尖点表示生成的块之和。

更一般地，对驯顺 $\ell$-分歧的 $G$，与 $\phi_\ell=1$ 相联系的因子 $\mathrm{Rep}_{\overline\Z_\ell,1}(G(F))$ 通常称为「幂幺 $\ell$-因子」。尽管如上例所示它可能不可分解，该模式所讲的是 $\mathrm{Rep}_{\Z_\ell}G(F)$ 如何分解为全子范畴之积，其中每个因子都等价于某个辅助的驯顺 $\ell$-分歧群的幂幺 $\ell$-因子。

以下是目前已知的：范畴 $\mathrm{Rep}_{\Z_\ell}(G(F))$ 的预期分解 (1) 由 [FS21] 的定理 IX.5.2 推出，其中 Fargues 与 Scholze 构造了从 $Z^1(W^0_F,\hatG)_{\Z_\ell}$ 上 $\hatG$-不变函数环到范畴 $\mathrm{Rep}_{\Z_\ell}(G(F))$ 的中心的一个映射。此外，他们的猜想 X.1.4 连同分解 (4.4) 必定蕴含所预期函子的存在（至少在导出层面），而它们与 Langlands 函子性的相容性将是「内建的」。

在这些发展之前，预期的分解 (1) 与 (3) 已在深度 $0$ 由 Lanard [Lan18] [Lan21] 相当显式地构造出来，并与所有已知的对应（DeBacker–Reeder，或经典群情形的 Arthur）相容。预期的范畴等价对 $G=\mathrm{GL}_n$（仍在深度 $0$）也已在 [Dat18] 中构造出来，尽管它们与 LLC 的相容性未知。
